\section{Method}
\label{sec:method}

\subsection{Estimating human preferences (\dfsp)}
\label{sec:dfsp}

We write the human preference probability as the judge's win probability plus a correction,
\[
  p_H(z)=p_J(z)+\delta(z),
\]
where $p_J(z)$ is the first coordinate of $z$. Instead of estimating $p_H$ from the human labels
alone, we estimate the correction $\delta$. When the judge is reasonably accurate, $\delta$ is small
or varies slowly, so it can be estimated from fewer labels than $p_H$.

\paragraph{Truncation.}
Following \citet{li2026personalizing}, fix a point $z$ and a parameter $\theta=(\theta_1,\theta_2)$
with $\theta_1\ge0$ and $0<\theta_2\le1$. Define
\[
  \omega_{\theta,z}(z')=p_J(z)+\min\bigl\{|p_J(z')-p_J(z)|,\ \theta_1\norm{z'-z}^{\theta_2}\bigr\}\,
  \operatorname{sgn}\{p_J(z')-p_J(z)\}.
\]
This limits how far the judge's prediction may move away from its value at $z$ as $z'$ moves away
from $z$. The parameter $\theta_1$ controls how much of the judge is used. For $\theta_1=0$,
$\omega_{\theta,z}$ is constant and the judge plays no role in the local fit. For large $\theta_1$,
$\omega_{\theta,z}=p_J$.

\paragraph{Estimator.}
Let $K_h(z',z)=\ind\{\norm{z'-z}_\infty\le h\}$. The estimated correction at $z$ is the local average
of the residuals $Y_i-\omega_{\theta,z}(Z_i)$,
\[
  \hat\delta_{\theta,h}(z)=\frac{\sum_{i\in\cI_n}K_h(Z_i,z)\{Y_i-\omega_{\theta,z}(Z_i)\}}
                                 {1\vee\sum_{i\in\cI_n}K_h(Z_i,z)},
\]
and the estimated preference probability is
\[
  \hat p_{\theta,h}(z)=\min\bigl[1,\max\{0,\ p_J(z)+\hat\delta_{\theta,h}(z)\}\bigr].
\]

\paragraph{Procedure.}
\begin{enumerate}
  \item Split the $n$ labeled comparisons into a fitting part and a validation part.
  \item For each $(\theta,h)$ in a grid that includes $\theta_1=0$, compute $\hat p_{\theta,h}$ on
        the fitting part.
  \item Choose $(\hat\theta,\hat h)$ with the smallest validation Brier score.
  \item Report $\hat p=\hat p_{\hat\theta,\hat h}$ on all comparisons.
\end{enumerate}
We call this estimator distributional few-shot personalization (\dfsp). The same procedure with $Z$
replaced by $p_J$ alone is the scalar version, \sfsp.

\subsection{Choosing which comparisons to label}
\label{sec:selection}

When the human budget is fixed, the choice of which comparisons to label affects the error. Let
$q$ be the density of $Z$ over all comparisons and $p$ the density of the labeled ones. Write
$\sigma^2(z)=p_H(z)\{1-p_H(z)\}$ for the label noise and $\gamma_1(z)$ for the local H\"older
constant of the correction $\delta$ near $z$, with H\"older exponent $\gamma_2$.

With a bandwidth chosen separately at each point, the error of \dfsp{} at $z$ is bounded by a bias
term $\gamma_1(z)^2h^{2\gamma_2}$ plus a variance term $\sigma^2(z)/\{np(z)h^d\}$. Minimizing the
$q$-weighted sum of these bounds over $p$ gives
\[
  p^{*}(z)\;\propto\;q(z)^{\frac{2\gamma_2+d}{4\gamma_2+d}}\;
  \gamma_1(z)^{\frac{2d}{4\gamma_2+d}}\;\sigma(z)^{\frac{4\gamma_2}{4\gamma_2+d}} .
\]
The rule assigns more labels to comparisons that are frequent ($q$ large), on which people
disagree ($\sigma$ large), and where the correction changes quickly ($\gamma_1$ large). A constant
correction is learned from few labels even if it is large. The rule therefore differs from
labeling the comparisons on which the judge is most uncertain.

The rule depends on unknown quantities, so we use it in two stages:
\begin{enumerate}
  \item Label a small random pilot sample and fit \dfsp{} on it.
  \item Estimate $\sigma$ from the pilot fit, and $\gamma_1$ from local differences of the fitted
        correction.
  \item Plug these estimates into $p^{*}$, mix the result with $q$, and draw the remaining labels
        from the pool with this density.
\end{enumerate}

\subsection{Post-training}
\label{sec:post-training}

Let $\pi_{\mathrm{ref}}$ be the initial model and $\pi_\vartheta$ the model being trained. As in DPO
\citep{rafailov2023dpo}, define
\[
  d_\vartheta(X)=\beta\Bigl[\log\frac{\pi_\vartheta(y_A\mid x)}{\pi_{\mathrm{ref}}(y_A\mid x)}
  -\log\frac{\pi_\vartheta(y_B\mid x)}{\pi_{\mathrm{ref}}(y_B\mid x)}\Bigr],
\]
so that $\sigma(d_\vartheta)$ is the model's implied probability that $y_A$ is preferred. We train on
all $N$ comparisons, using the estimated probabilities $\hat p_i=\hat p(Z_i)$ as targets:
\[
  \widehat R(\vartheta;\hat p)=-\frac1N\sum_{i=1}^{N}
  \Bigl[\hat p_i\log\sigma\{d_\vartheta(X_i)\}+(1-\hat p_i)\log\sigma\{-d_\vartheta(X_i)\}\Bigr].
\]
This is the usual soft-label DPO loss \citep{lee2024rlaif,furuta2024gapo}. The difference from
existing work is that the targets come from \dfsp{} rather than directly from the judge.
